\documentclass[a4paper,12pt]{article}
\usepackage{fontspec}
\usepackage{polyglossia}
\setmainlanguage{russian}
\setotherlanguage{english}
\setmainfont{Times New Roman}
\newfontfamily\cyrillicfonttt{Consolas}
\usepackage{amsmath,amssymb,amsthm}
\usepackage{geometry}
\usepackage{hyperref}
\usepackage{enumitem}
\usepackage{xcolor}
\usepackage{graphicx}
\usepackage{indentfirst}
\usepackage{array}
\geometry{left=3cm,right=1.5cm,top=2cm,bottom=2cm}
\hypersetup{unicode=true,colorlinks=true,linkcolor=blue,urlcolor=blue}
\setlist[enumerate]{itemsep=0.2em,topsep=0.4em}
\setlength{\parindent}{1.25cm}
\setlength{\parskip}{0.35em}
\linespread{1.5}
\setcounter{secnumdepth}{0}
\newtheorem{assumption}{Предположение}
\newtheorem{proposition}{Утверждение}
\newtheorem{remark}{Замечание}
\newcommand{\thesistitle}{Численное решение терминальной задачи оптимального управления с линейной динамикой}
\newcommand{\makevkrtitle}{%
  \begin{titlepage}
  \thispagestyle{empty}
  \begin{center}
  \linespread{1}\selectfont
  \includegraphics[width=8cm,height=4cm,keepaspectratio]{论文要求/论文模板/logo.jpg}\par
  \vspace{0.25cm}
  Московский государственный университет имени М.\,В.~Ломоносова\\
  \vspace{0.1cm}
  Факультет вычислительной математики и кибернетики\\
  \vspace{0.1cm}
  Кафедра оптимального управления

  \vspace{3cm}
  {\Large Сяо Ижу\par}

  \vspace{1cm}
  {\bfseries\LARGE \thesistitle\par}
  \vspace{2cm}
  МАГИСТЕРСКАЯ ДИССЕРТАЦИЯ
  \end{center}

  \vspace{2cm}
  \begin{flushright}
  {\bfseries Научный руководитель:}\\
  старший преподаватель\\
  Кулевский Александр Владиславович
  \end{flushright}

  \vfill
  \centerline{Москва, 2026}
  \end{titlepage}
}

\begin{document}
\makevkrtitle
\setcounter{page}{2}
\tableofcontents
\newpage

\section{Аннотация}

В работе рассматриваются терминальные задачи оптимального управления для линейной системы
\[
\dot{x}(t)=Ax(t)+Bu(t),\quad x(0)=x_0,\quad u(t)\in U,\quad t\in[0,T],
\]
где $U$ --- компактное выпуклое множество управлений. Исследуются две постановки: минимизация квадратичной терминальной ошибки и минимизация линейного терминального функционала при аффинном терминальном ограничении. Основная идея состоит в замене поиска по всей функции управления конечномерным поиском по терминальному сопряжённому параметру $q=\psi(T)$ или по множителю $\lambda$.

Численная схема строится как комбинация непрерывных необходимых условий принципа максимума и дискретной рекурсии Euler для состояния. Поэтому программа решает не полный дискретный KKT-комплекс, а задачу минимизации дискретной параметрической невязки, индуцированной непрерывными условиями. В численных экспериментах сравниваются направления Гаусса--Ньютона и градиента, влияние сглаживания в bang-bang постановке, выигрыш по размерности в четырёхмерной системе, а также поведение метода в задаче с терминальным ограничением. Результаты показывают, что низкоразмерная параметризация вычислима и удобна для интерпретации, но требует осторожности при негладкости, вырождении и повышенных требованиях к точности.

\section{Ключевые слова}

оптимальное управление; линейная динамика; терминальная задача; численное решение; принцип максимума Понтрягина; сопряжённое уравнение; опорная функция; конечномерная параметризация; метод Гаусса--Ньютона

\section{Глава 1. Введение}

\subsection{1.1 Актуальность и значение исследования}

Теория оптимального управления изучает, как при заданной динамической системе, ограничениях и критерии качества выбирать управляющую функцию так, чтобы движение системы было оптимальным в выбранном смысле\cite[с. 13--27]{pontryagin1969}\cite[с. 76--91]{lee1972}\cite[с. 8--15]{blagodatskikh2001}. Терминальные задачи оптимального управления ориентированы на свойства состояния в конечный момент времени: требуется либо максимально приблизить конечное состояние к заданной цели, либо обеспечить его принадлежность заданному целевому множеству и одновременно оптимизировать терминальный функционал.

Для линейных систем с постоянными коэффициентами уравнение состояния, матричная экспонента, множество достижимости, опорная функция и сопряжённое уравнение обладают достаточно прозрачной структурой. Поэтому линейные системы являются не только базовым объектом теории оптимального управления, но и естественной основой для построения вычислимых численных методов.

Цель настоящей работы состоит в том, чтобы использовать эту структуру для перехода от бесконечномерной оптимизации по всей функции управления к конечномерному поиску по терминальному сопряжённому параметру или по множителю\cite[с. 180--190]{kiselev2007}\cite[с. 5--7]{avvakumov1995}. В отличие от прямых методов, где независимыми переменными являются значения управления во всех узлах сетки, здесь размерность численно оптимизируемой переменной обычно связана лишь с размерностью пространства состояний или с числом терминальных ограничений.

В работе рассматривается только класс линейных терминальных задач с наиболее прозрачной структурой. Вне рассмотрения остаются нелинейные системы, фазовые ограничения, общая теория глобальной сходимости негладких методов и полная теория дискретных оптимальностей для разностных моделей. Основной акцент сделан на том, как из классических необходимых условий построить вычислимую систему низкой размерности и как корректно интерпретировать её решения.

\subsection{1.2 Современное состояние исследований}

Для рассматриваемого класса задач ключевую роль играют три группы результатов. Во-первых, принцип максимума Понтрягина задаёт сопряжённое уравнение, условие максимума и трансверсальные условия\cite[с. 23--27]{pontryagin1969}\cite[с. 336--371]{lee1972}\cite{liberzon2012}. Во-вторых, для линейных систем конечное состояние зависит от управления аффинно\cite[с. 15--59]{kiselev2007}. Множество достижимости и опорная функция дают удобный геометрический язык для описания структуры управления\cite[с. 29--37]{blagodatskikh2001}.

В-третьих, численные методы обычно делятся на прямые и косвенные\cite{betts2010,rao2009,bock1984,hager2000,dontchev2000}. Прямые методы оптимизируют дискретизированное управление, тогда как косвенные сначала используют необходимые условия. В настоящей работе используется косвенная идея, но вычислительный объект выбирается в более компактной форме: вместо полной краевой задачи решается конечномерная задача по терминальному параметру, а результаты интерпретируются через дискретную невязку, терминальную ошибку и невязку ограничения.

\subsection{1.3 Содержание работы и структура исследования}

Работа посвящена двум линейным терминальным постановкам: задаче минимизации квадратичной терминальной ошибки и задаче минимизации линейного терминального функционала при аффинном ограничении на конечное состояние. Основные подзадачи таковы:
\begin{enumerate}
\item зафиксировать линейную модель, терминальные критерии и параметризацию по $q$ или $\lambda$;
\item вывести необходимые условия и соответствующие уравнения дискретной невязки;
\item построить численный метод на основе Euler, конечных разностей, Армихо и Гаусса--Ньютона;
\item сопоставить полученные кандидатные решения с прямым дискретным ориентиром и описать границы интерпретации результатов.
\end{enumerate}

Структура работы такова: во второй главе излагаются математические основы линейной системы и терминальных задач; в третьей главе выводятся условия оптимальности и конечномерная параметризация; в четвёртой главе строится численный метод; в пятой главе приводятся численные эксперименты и обсуждение результатов; в шестой главе формулируется итоговая позиция метода.

\section{Глава 2. Математические основы линейной системы управления и терминальных задач оптимального управления}

\subsection{2.1 Линейная система управления}

Рассмотрим линейную систему с постоянными коэффициентами\cite[с. 15--31]{kiselev2007}\cite[с. 76--91]{lee1972}
\[
\dot{x}(t)=Ax(t)+Bu(t),\quad t\in[0,T],
\]
\[
x(0)=x_0.
\]
Здесь $x(t)\in\mathbb{R}^n$, $u(t)\in\mathbb{R}^m$, $A\in\mathbb{R}^{n\times n}$ и $B\in\mathbb{R}^{n\times m}$.

\begin{assumption}[Основные предположения]
Всюду далее предполагается, что $T>0$, матрицы $A$ и $B$ заданы, множество $U\subset\mathbb{R}^m$ непусто, компактно и выпукло, а допустимые управления принадлежат пространству $L_\infty([0,T];U)$. Для задач с терминальными ограничениями основным случаем считается аффинное множество
\[
M=\{x\in\mathbb{R}^n:Cx=d\}.
\]
\end{assumption}

\subsection{2.2 Матричная экспонента и формула решения системы}

Для однородной системы $\dot{x}=Ax$ решение задаётся матричной экспонентой $e^{At}$\cite[с. 15--31]{kiselev2007}\cite[с. 76--91]{lee1972}. По формуле вариации постоянных решение управляемой системы имеет вид
\[
x(t)=e^{At}x_0+\int_0^t e^{A(t-s)}Bu(s)\,ds.
\]
В частности,
\[
x(T)=e^{AT}x_0+\int_0^T e^{A(T-s)}Bu(s)\,ds.
\]
Из этой формулы видно, что конечное состояние зависит от функции управления аффинно.

\subsection{2.3 Допустимые управления, множества достижимости и опорная функция}

Положим
\[
\mathcal{U}=L_\infty([0,T];U),\qquad
X(T)=\{x(T;u):u(\cdot)\in\mathcal U\}.
\]
Из формулы решения следует, что терминальное отображение аффинно по управлению, а при выпуклом $U$ множество достижимости $X(T)$ также выпукло\cite[с. 76--91]{lee1972}\cite[с. 32--40]{kiselev2007}. Для численной части далее используется дискретная рекурсия Euler
\[
y_{m+1}=y_m+\tau(Ay_m+Bu_m),\qquad \tau=T/M,
\]
которая задаёт дискретное терминальное отображение и соответствующее дискретное множество достижимости $X_M(T)$.

Для непустого компактного выпуклого множества $U\subset\mathbb{R}^m$ вводится опорная функция
\[
c_U(\eta)=\max_{u\in U}\langle \eta,u\rangle.
\]
Именно она позволяет переписать условие максимума как линейную максимизацию по $U$ и затем восстановить управление из величины $B^\top\psi(t)$. Для ящичных и многогранных ограничений эта функция обычно негладка, что важно для дальнейшей численной интерпретации.

\section{Глава 3. Условия оптимальности для линейных терминальных задач оптимального управления}

\subsection{3.1 Единая минимизационная запись задачи}

Все исходные задачи далее записываются как задачи минимизации
\[
J(u)\to\min.
\]
Для общей терминальной постановки имеем
\[
\dot{x}(t)=Ax(t)+Bu(t),\quad x(0)=x_0,\quad u(t)\in U,
\]
\[
J(u)=\Phi(x(T))\to\min,
\]
а при наличии терминального ограничения дополнительно
\[
x(T)\in M.
\]

Чтобы использовать опорную функцию в максимизационной форме, введём сопряжённую переменную
\[
\psi(t)=-p(t),
\]
где $p(t)$ --- стандартная сопряжённая переменная минимизационной записи.

\subsection{3.2 Гамильтониан и сопряжённое уравнение}

При использовании максимизационной формы сопряжённой переменной гамильтониан определяется как
\[
H(x,u,\psi,t)=\langle \psi,Ax+Bu\rangle.
\]
Так как система линейна по состоянию, сопряжённое уравнение имеет вид
\[
\dot{\psi}(t)=-A^\top\psi(t).
\]
Если положить
\[
q=\psi(T),
\]
то решение сопряжённого уравнения записывается как
\[
\psi(t,q)=e^{A^\top(T-t)}q.
\]
Следовательно, для линейной динамики вся сопряжённая траектория полностью определяется конечномерным терминальным вектором $q$.

\subsection{3.3 Условие максимума и представление через опорную функцию}

Принцип максимума Понтрягина даёт условие\cite[с. 23--27]{pontryagin1969}\cite[с. 336--371]{lee1972}\cite[с. 8--15]{blagodatskikh2001}
\[
H(x^*(t),u^*(t),\psi(t),t)=\max_{u\in U}H(x^*(t),u,\psi(t),t).
\]
Поскольку управление входит в гамильтониан линейно, это условие эквивалентно записи
\[
u^*(t)\in\arg\max_{u\in U}\langle B^\top\psi(t),u\rangle.
\]
С использованием опорной функции получаем
\[
\langle B^\top\psi(t),u^*(t)\rangle=c_U(B^\top\psi(t)).
\]
Если $c_U$ дифференцируема, то
\[
u^*(t)=\nabla c_U(B^\top\psi(t)).
\]

\subsection{3.4 Терминальное условие для первой задачи}

Для первой задачи терминальный функционал имеет вид
\[
J_1(u)=\frac{1}{2}\|x(T;u)-x_T\|^2.
\]
Положим
\[
\Phi_1(x)=\frac{1}{2}\|x-x_T\|^2.
\]
Тогда
\[
\nabla\Phi_1(x)=x-x_T.
\]
Для стандартной минимизационной сопряжённой переменной получаем терминальное условие
\[
p(T)=x^*(T)-x_T.
\]
Так как в работе используется соглашение $\psi=-p$, имеем
\[
\psi(T)=x_T-x^*(T).
\]
Если положить $q=\psi(T)$, то для оптимального решения должно выполняться
\[
q=x_T-x(T,q).
\]

\begin{remark}
Если целевая точка $x_T$ достижима и оптимальная терминальная ошибка равна нулю, то получается $q=0$. В этом случае условие максимума перестаёт однозначно определять управление. Поэтому параметризация по $q$ особенно информативна в задачах с ненулевым терминальным отклонением, на границе достижимости или после сглаживания, но требует дополнительной интерпретации в вырожденном случае.
\end{remark}

\subsection{3.5 Замечания об интерпретации, вырождении и негладкости}

Для первой задачи функционал $J_1(u)=\frac12\|x(T;u)-x_T\|^2$ выпукл по управлению, однако численная реализация в работе используется не как доказательство строгого глобального оптимума, а как вычислимое приближение необходимых условий. Поэтому даже в выпуклой постановке нужно различать непрерывную задачу, её дискретную реализацию и фактически полученную терминальную ошибку.

Если цель достижима и оптимальная терминальная ошибка равна нулю, то возникает вырождение $q=0$, при котором условие максимума уже не задаёт управление однозначно. Для ящичных ограничений дополнительную трудность создаёт негладкость опорной функции: при $\eta_i(t)=0$ максимизация по соответствующей компоненте становится многозначной, а соглашение $\operatorname{sign}(0)=0$ остаётся лишь численным правилом. Именно поэтому далее используется сглаженная аппроксимация управления и отдельно отслеживаются величина невязки, терминальная ошибка и параметр сглаживания.

\subsection{3.6 Трансверсальное условие для задачи с терминальным целевым множеством}

Для второй задачи
\[
J_2(u)=a^\top x(T;u)\to\min,\quad x(T)\in M
\]
при общем выпуклом терминальном множестве имеем\cite[с. 121--127]{pontryagin1969}\cite[с. 336--371]{lee1972}
\[
p(T)\in a+N_M(x^*(T)),
\]
где $N_M(x^*(T))$ --- нормальный конус к множеству $M$ в точке $x^*(T)$. Поскольку $\psi=-p$, получаем
\[
\psi(T)\in -a-N_M(x^*(T)).
\]

Если
\[
M=\{x\in\mathbb{R}^n:Cx=d\},
\]
то нормальное пространство совпадает с $\operatorname{Range}(C^\top)$. Если, как в основной части работы, матрица $C\in\mathbb{R}^{r\times n}$ имеет полный построчный ранг, то существует множитель $\lambda\in\mathbb{R}^r$ такой, что
\[
p(T)=a+C^\top\lambda,
\]
и потому
\[
\psi(T)=-a-C^\top\lambda.
\]

\subsection{3.7 Идея конечномерной параметризации}

Вместо поиска по всей функции управления используется конечномерная параметризация. Для первой задачи свободным параметром служит $q=\psi(T)\in\mathbb{R}^n$, а для второй --- множитель $\lambda$, задающий терминальный вектор
\[
q(\lambda)=-a-C^\top\lambda.
\]
После этого по формулам сопряжённого уравнения и условия максимума восстанавливаются управление, траектория состояния и терминальная невязка. Тем самым бесконечномерная задача по управлению заменяется конечномерным поиском по $q$ или $\lambda$\cite[с. 180--190]{kiselev2007}\cite[с. 5--7]{avvakumov1995}.

\section{Глава 4. Численный метод на основе принципа максимума}

\subsection{4.1 Общая идея численного решения}

В работе используется косвенный подход\cite[с. 23--27]{pontryagin1969}\cite[с. 3--7]{avvakumov1995}. Переменной численного поиска служит не вся функция управления на сетке, а конечномерный параметр $q$ или $\lambda$. Вычислительная цепочка имеет вид
\[
z\longmapsto q(z)\longmapsto \psi(t,z)\longmapsto u(t,z)\longmapsto x(t,z)\longmapsto R_M(z),
\]
где $z=q$ для первой задачи и $z=\lambda$ для второй, а $R_M$ --- дискретная параметрическая невязка на сетке из $M$ шагов. Тем самым программа реализует схему ``непрерывная терминальная параметризация + дискретная рекурсия Euler'', а итог интерпретируется как дискретное кандидатное решение, согласованное с непрерывными необходимыми условиями.

\subsection{4.2 Временная дискретизация и вычислительное отображение}

Разобьём отрезок $[0,T]$ на $M$ равных частей:
\[
t_m=m\tau,\quad \tau=\frac{T}{M},\quad m=0,1,\ldots,M.
\]
Для состояния применяется явная схема Euler:
\[
x_{m+1}=x_m+\tau(Ax_m+Bu_m),\quad x_0=x_0.
\]
Эквивалентно,
\[
x_{m+1}=A_dx_m+B_du_m,\quad A_d=I+\tau A,\quad B_d=\tau B.
\]

При заданном терминальном параметре $q$ дискретные значения сопряжённой переменной берутся в виде
\[
\psi_m(q)=e^{A^\top(T-t_m)}q.
\]
Положим
\[
\eta_m(q)=B^\top\psi_m(q).
\]
Тогда управление задаётся формулой
\[
u_m(q)\in\arg\max_{u\in U}\langle \eta_m(q),u\rangle.
\]
Если опорная функция дифференцируема, то
\[
u_m(q)=\nabla c_U(\eta_m(q)).
\]
После этого по рекурсии вычисляется $x_M(q)$.

Таким образом, при фиксированном $q$ или $\lambda$ весь численный процесс определяется однозначно: сначала строится сопряжённая траектория, затем управление, после чего по дискретной рекурсии получается траектория состояния и терминальная невязка.

\subsection{4.3 Уравнение невязки для задачи квадратичной терминальной ошибки}

Для первой задачи в качестве численно оптимизируемой переменной берётся
\[
z=q\in\mathbb{R}^n.
\]
По параметру $q$ строится управление $u_m(q)$, затем терминальное состояние $x_M(q)$. Согласно непрерывному терминальному условию
\[
q=x_T-x(T,q),
\]
на дискретном уровне вводится приближённая терминальная невязка
\[
F(q)=q+x_M(q)-x_T.
\]
Соответствующая конечномерная численная задача имеет вид
\[
\min_{q\in\mathbb{R}^n}\Theta_1(q),\quad
\Theta_1(q)=\frac{1}{2}\|F(q)\|^2.
\]

После нахождения параметра $q^*$ строятся последовательности
\[
u_m^*=u_m(q^*),\quad x_m^*=x_m(q^*),\quad m=0,1,\ldots,M.
\]
При этом обязательно сообщаются как величина терминальной приближённой невязки,
\[
\|F(q^*)\|,
\]
так и фактическая терминальная ошибка
\[
e_T=\|x_M(q^*)-x_T\|.
\]
Если цель недостижима, величина $e_T$ не обязана стремиться к нулю; если цель достижима, но возникает вырождение $q=0$, то одна лишь невязка может быть недостаточна для однозначного восстановления оптимального управления.

\subsection{4.4 Невязка для задачи с терминальным целевым множеством}

Во второй задаче рассматривается аффинное терминальное множество
\[
M=\{x\in\mathbb{R}^n:Cx=d\}.
\]
Тогда терминальный сопряжённый параметр выражается через множитель:
\[
q(\lambda)=-a-C^\top\lambda,\quad \lambda\in\mathbb{R}^r,
\]
где предполагается, что $C\in\mathbb{R}^{r\times n}$ имеет полный построчный ранг.

При фиксированном $\lambda$ вычисляются
\[
\psi_m(\lambda)=e^{A^\top(T-t_m)}q(\lambda),\quad
u_m(\lambda)\in\arg\max_{u\in U}\langle B^\top\psi_m(\lambda),u\rangle,
\]
после чего по разностной схеме строится $x_M(\lambda)$. Затем вводится терминальная невязка ограничения
\[
G(\lambda)=Cx_M(\lambda)-d.
\]
Численная задача формулируется как
\[
\min_{\lambda\in\mathbb{R}^r}\Theta_2(\lambda),\quad
\Theta_2(\lambda)=\frac{1}{2}\|G(\lambda)\|^2.
\]

Здесь важно, что функция $\Theta_2(\lambda)$ не является прямым переписыванием исходного терминального функционала как численного критерия по $\lambda$. Она служит численным способом найти параметр, который даёт приблизительно допустимое кандидатное решение, согласованное с непрерывными необходимыми условиями.

После нахождения $\lambda^*$ сообщаются
\[
r_T=\|G(\lambda^*)\|,\quad J_2=a^\top x_M(\lambda^*).
\]
Если $r_T$ остаётся выше принятой точности, это означает лишь неполное выполнение терминального ограничения в рамках выбранной дискретной процедуры, а не строгую оптимальность.

\subsection{4.5 Якобиан конечных разностей и направление Гаусса--Ньютона}

Обозначим через $R(z)$ общую невязку: $R=F$ в первой задаче и $R=G$ во второй. Пусть
\[
R(z)\in\mathbb{R}^r,\qquad z\in\mathbb{R}^s.
\]
Тогда минимизируется функция
\[
\Theta(z)=\frac{1}{2}\|R(z)\|^2.
\]

Чтобы не выводить вариационное уравнение для каждой задачи, Jacobian аппроксимируется прямыми конечными разностями. Если $e_j$ --- $j$-й базисный вектор, то
\[
J_R(z)_{:,j}\approx \frac{R(z+he_j)-R(z)}{h},\quad j=1,\ldots,s.
\]
Соответственно,
\[
\nabla\Theta(z)\approx J_R(z)^\top R(z).
\]

Градиентное направление берётся как
\[
d_k=-J_R(z_k)^\top R(z_k).
\]
Для направления Гаусса--Ньютона на шаге $k$ линеаризуется невязка:
\[
R(z_k+d)\approx R(z_k)+J_R(z_k)d.
\]
Это приводит к линейному наименьшему квадрату
\[
\min_d \frac12\|R(z_k)+J_R(z_k)d\|^2,
\]
нормальные уравнения которого имеют вид
\[
J_R(z_k)^\top J_R(z_k)d_k=-J_R(z_k)^\top R(z_k).
\]
Для повышения устойчивости используется демпфированная форма
\[
\left(J_R(z_k)^\top J_R(z_k)+\alpha I\right)d_k=-J_R(z_k)^\top R(z_k),
\]
где $\alpha\ge 0$ --- параметр регуляризации.

\subsection{4.6 Линейный поиск Армихо и критерии остановки}

После выбора направления $d_k$ шаг определяется линейным поиском Армихо. При начальном значении $\rho=1$ и параметрах $\beta,\sigma\in(0,1)$ шаг уменьшается до тех пор, пока не выполнится условие
\[
\Theta(z_k+\rho d_k)\le
\Theta(z_k)+\sigma\rho\nabla\Theta(z_k)^\top d_k.
\]
Затем выполняется обновление
\[
z_{k+1}=z_k+\rho d_k.
\]

Остановка итераций происходит по одному из критериев:
\[
\|R(z_k)\|\le \varepsilon_R,
\]
или
\[
\|\nabla\Theta(z_k)\|\le \varepsilon_g,
\]
или по достижении максимального числа итераций $K_{\max}$. Если достигнут предел $K_{\max}$ без выполнения основного критерия по невязке, в отчёте фиксируется текущая лучшая точка и отдельно указывается, что целевая точность по невязке не была достигнута.

\subsection{4.7 Негладкие множества управлений и сглаживание}

Если множество управлений является интервалом
\[
U=[-u_{\max},u_{\max}],
\]
то
\[
c_U(\eta)=u_{\max}|\eta|,\quad
u_m=u_{\max}\operatorname{sign}(\eta_m).
\]
Для многомерного ящика
\[
U=\{u\in\mathbb{R}^m: |u_i|\le u_{\max,i},\ i=1,\ldots,m\}
\]
получаем
\[
c_U(\eta)=\sum_{i=1}^m u_{\max,i}|\eta_i|,\quad
u_{m,i}=u_{\max,i}\operatorname{sign}(\eta_{m,i}).
\]

Такая структура соответствует bang-bang управлению и создаёт негладкость на переключениях. Для повышения устойчивости вводится сглаженная опорная функция\cite[с. 37--51]{blagodatskikh2001}\cite[с. 8--9]{avvakumov1995}
\[
c_{U,\mu}(\eta)=
\sum_{i=1}^m u_{\max,i}\sqrt{\eta_i^2+\mu^2},
\quad \mu>0.
\]
Тогда управление принимает вид
\[
u_{m,i}^{\mu}=
\frac{u_{\max,i}\eta_{m,i}}{\sqrt{\eta_{m,i}^2+\mu^2}}.
\]
При $\mu\to 0$ оно приближается к bang-bang форме.

Нужно подчеркнуть, что при $\mu>0$ решается сглаженная аппроксимация, а не исходная негладкая задача. Чем меньше $\mu$, тем ближе управление к исходной bang-bang структуре, но тем чувствительнее численная невязка к переключениям и тем менее надёжны конечные разности вблизи точек смены режима.

\subsection{4.8 Условия применимости и признак низкой размерности}

Поведение метода определяется тремя факторами: гладкостью невязки, обусловленностью матрицы $J_k^\top J_k$ и размерностью оптимизируемой переменной. При $\mu>0$ сглаживание обычно повышает устойчивость конечных разностей и шагов Гаусса--Ньютона, тогда как негладкость при $\mu=0$ усиливает чувствительность к переключениям. При плохой обусловленности демпфирование повышает устойчивость шага, но не устраняет самой причины вырождения.

Главное структурное преимущество схемы состоит в низкой размерности: в первой задаче оптимизируется параметр размерности $n$, во второй --- множитель размерности $r$, тогда как прямой дискретный метод работает с размерностью $M\cdot m$. Поэтому сильная сторона подхода --- компактная параметризация и интерпретируемость, а не универсальная гарантия лучшей точности.

\subsection{4.9 Алгоритм 1: задача квадратичной терминальной ошибки}

Входные данные:
\[
A,\ B,\ x_0,\ x_T,\ U,\ T,\ M,\ q_0,\ h,\ \varepsilon_R,\ \varepsilon_g,\ K_{\max}.
\]
Алгоритм состоит из следующих шагов:
\begin{enumerate}
\item положить $\tau=T/M$, задать начальное значение $q_0$ и номер итерации $k=0$;
\item по текущему $q_k$ вычислить $\psi_m(q_k)$, $\eta_m(q_k)$ и $u_m(q_k)$;
\item по схеме Euler вычислить траекторию состояния и терминальное значение $x_M(q_k)$;
\item вычислить невязку $F(q_k)=q_k+x_M(q_k)-x_T$ и значение $\Theta_1(q_k)=\frac12\|F(q_k)\|^2$;
\item если $\|F(q_k)\|\le\varepsilon_R$, остановиться;
\item вычислить конечнодифференциальный Jacobian и выбрать направление градиентного или гаусс-ньютоновского типа;
\item выполнить поиск Армихо и обновить параметр $q_{k+1}=q_k+\rho_k d_k$;
\item если выполнен критерий по градиенту или превышено $K_{\max}$, остановиться; иначе перейти к следующей итерации.
\end{enumerate}
Выходные данные:
\[
q^*,\quad \{u_m^*\}_{m=0}^{M-1},\quad \{x_m^*\}_{m=0}^{M},\quad
\|F(q^*)\|,\quad e_T=\|x_M^*-x_T\|.
\]

\subsection{4.10 Алгоритм 2: задача с аффинным терминальным целевым множеством}

Входные данные:
\[
A,\ B,\ x_0,\ U,\ T,\ M,\ a,\ C,\ d,\ \lambda_0,\ h,\ \varepsilon_R,\ \varepsilon_g,\ K_{\max}.
\]
Шаги алгоритма:
\begin{enumerate}
\item положить $\tau=T/M$, задать начальное значение $\lambda_0$ и номер итерации $k=0$;
\item вычислить терминальный сопряжённый параметр $q(\lambda_k)=-a-C^\top\lambda_k$;
\item построить $\psi_m(\lambda_k)$, $\eta_m(\lambda_k)$ и управление $u_m(\lambda_k)$;
\item по дискретной системе получить траекторию состояния и значение $x_M(\lambda_k)$;
\item вычислить невязку $G(\lambda_k)=Cx_M(\lambda_k)-d$ и целевую функцию $\Theta_2(\lambda_k)=\frac12\|G(\lambda_k)\|^2$;
\item если $\|G(\lambda_k)\|\le\varepsilon_R$, остановиться;
\item вычислить якобиан, определить направление и выполнить поиск Армихо;
\item обновить множитель $\lambda_{k+1}=\lambda_k+\rho_k d_k$;
\item остановиться по критерию градиента или по пределу числа итераций.
\end{enumerate}
Выход:
\[
\lambda^*,\quad q(\lambda^*),\quad \{u_m^*\}_{m=0}^{M-1},\quad
\{x_m^*\}_{m=0}^{M},\quad r_T=\|G(\lambda^*)\|,\quad J_2=a^\top x_M^*.
\]

\section{Глава 5. Численные эксперименты и анализ результатов}

Использованные в работе программы, таблицы результатов и графические материалы размещены в открытом GitHub-репозитории\cite{github-code}, что позволяет воспроизводить вычисления и проверять детали реализации.

\subsection{5.1 Экспериментальная настройка и показатели качества}

В этой главе численно проверяются вычислимость и устойчивость алгоритмов главы 4. Если не оговорено иное, используется явная схема Euler, а общие параметры расчёта сведены в таблицу.
\[
\begin{array}{c|c}
\text{параметр} & \text{значение}\\
\hline
M & 50,100,200,400,800\ \text{или значение, указанное в примере}\\
h & 10^{-6}\\
\varepsilon_R & 10^{-6}\\
\varepsilon_g & 10^{-8}\\
K_{\max} & 120\\
\beta,\sigma & 0.5,\ 10^{-4}\\
\alpha & 10^{-8}
\end{array}
\]

Во всех примерах фиксируются величины $\Theta_k=\frac12\|R(z_k)\|^2$, $g_k=\|J_R(z_k)^\top R(z_k)\|$ и $d_k=\|z_{k+1}-z_k\|$. Для первой задачи дополнительно используется терминальная ошибка $e_T=\|x_M-x_T\|$, а для второй --- невязка ограничения $r_T=\|Cx_M-d\|$ и значение функционала $J_2=a^\top x_M$. Ключевые выводы строятся именно по этим величинам.

\subsection{5.2 Два типа задач и распределение примеров}

Все численные примеры согласованы с двумя основными задачами работы и обозначены как case1--case5. Три примера относятся к задаче минимизации квадратичной терминальной ошибки, а два --- к задаче с линейным терминальным функционалом и аффинным терминальным ограничением. Для удобства их доказательная роль сведена в таблицу~\ref{tab:case-overview}.

\begin{table}[htbp]
\centering
\small
\caption{Обзор доказательной роли пяти основных примеров}
\label{tab:case-overview}
\begin{tabular}{|c|c|>{\raggedright\arraybackslash}p{3.4cm}|>{\raggedright\arraybackslash}p{7.0cm}|}
\hline
case & Тип задачи & Объекты сравнения & Поддерживаемый вывод \\
\hline
case1 & Первая задача & Гаусс--Ньютон, градиентный метод, прямой метод, изменение сетки & Показывает различие итерационного поведения, задаёт базу по точности и фиксирует слабое влияние сетки в этом примере. \\
\hline
case2 & Первая задача & Разные значения $\mu$ и bang-bang выбор & Показывает, что сглаживание одновременно влияет на устойчивость сходимости и форму управления. \\
\hline
case3 & Первая задача & Параметризация по сопряжённой переменной и прямой метод & Показывает сочетание меньшей размерности переменной и ограничения по достигаемой терминальной точности. \\
\hline
case4 & Вторая задача & Геометрия терминального множества, множитель, линейная цель & Показывает, как терминальное множество и линейная цель совместно определяют кандидатное конечное состояние. \\
\hline
case5 & Вторая задача & Разные начальные значения множителя & Показывает устойчивость невязки ограничения и итогового значения цели по отношению к начальному множителю. \\
\hline
\end{tabular}
\end{table}

\subsection{5.3 Примеры для первой задачи}

В этом разделе рассматривается первая задача $J_1(u)=\frac12\|x(T)-x_T\|^2\to\min$. Во всех трёх примерах используется терминальный параметр $q$, дискретная невязка $F(q)=q+x_M(q)-x_T$ и схема наименьших квадратов по этой невязке; при необходимости для контроля привлекается прямой дискретный метод.

\paragraph{case1: двумерная геометрия терминального отклонения и сравнение методов}
Рассматривается двумерная система
\[
\dot{x}=Ax+Bu,\quad |u|\le 1,
\]
где
\[
A=
\begin{pmatrix}
0&1\\
-1&0
\end{pmatrix},
\qquad
B=
\begin{pmatrix}
0\\
1
\end{pmatrix}.
\]
Выбираются параметры
\[
x_0=
\begin{pmatrix}
1\\
0
\end{pmatrix},
\quad
x_T=
\begin{pmatrix}
0.6\\
-0.6
\end{pmatrix},
\quad
T=2.
\]
Функционал имеет вид
\[
J_1(u)=\frac12\|x(T)-x_T\|^2.
\]
В примере сравниваются направление Гаусса--Ньютона, градиентное направление и прямой дискретный базовый метод.

\paragraph{case2: негладкое управление и сравнение параметров сглаживания}
Здесь сохраняются двумерная система и ограничение управления из case1, но сопоставляются негладкое bang-bang управление
\[
u_m=\operatorname{sign}(\eta_m),\quad \eta_m=B^\top\psi_m
\]
со сглаженным управлением
\[
u_m^\mu=\frac{\eta_m}{\sqrt{\eta_m^2+\mu^2}}.
\]
Берутся значения
\[
\mu\in\{10^{-1},10^{-2},10^{-3}\},
\]
а случай $\mu=0$ служит контрольным вариантом без сглаживания. Фиксируются
\[
M=200,\qquad q_0=
\begin{pmatrix}
0\\
0
\end{pmatrix}.
\]

\paragraph{case3: четырёхмерная линейная система средней размерности}
Чтобы показать эффект снижения размерности переменной в более высокой размерности, рассматривается система
\[
A=
\begin{pmatrix}
0&1&0&0\\
-1&-0.2&0.3&0\\
0&0&0&1\\
0.2&0&-1&-0.1
\end{pmatrix},
\qquad
B=
\begin{pmatrix}
0&0\\
1&0\\
0&0\\
0&1
\end{pmatrix}.
\]
Выбираются
\[
x_0=(1,0,-0.5,0)^\top,\quad x_T=(0,0,0,0)^\top,\quad T=3,
\]
при ограничениях
\[
|u_1|\le 1,\qquad |u_2|\le 1.
\]
Функционал первой задачи:
\[
J_1(u)=\frac12\|x(T)-x_T\|^2\to\min.
\]
Число временных шагов берётся равным $M=200$.

\subsection{5.4 Примеры для второй задачи: минимизация линейного терминального функционала при терминальном ограничении}

В этом разделе рассматривается задача $x(T)\in M,\ J_2(u)=a^\top x(T)\to\min$. Главной неизвестной здесь является множитель $\lambda$, связанный с терминальным ограничением.

\paragraph{case4: геометрия целевого множества и линейный терминальный функционал}
Рассматривается двумерная вращательная система
\[
A=
\begin{pmatrix}
0&1\\
-1&0
\end{pmatrix},
\qquad
B=
\begin{pmatrix}
0\\
1
\end{pmatrix},
\qquad |u|\leq 1.
\]
Параметры:
\[
x_0=
\begin{pmatrix}
0\\
1
\end{pmatrix},
\quad
T=2,
\quad
a=
\begin{pmatrix}
1\\
0
\end{pmatrix}.
\]
Терминальное множество задаётся как
\[
M=\{x\in\mathbb{R}^2:Cx=d\},
\quad
C=
\begin{pmatrix}
0&1
\end{pmatrix},
\quad
d=-0.2.
\]
То есть конечное состояние должно лежать на прямой $x_2=-0.2$, а целевая функция равна
\[
J_2(u)=a^\top x(T)\to\min,\qquad Cx(T)=d.
\]
По трансверсальному условию
\[
q(\lambda)=-a-C^\top\lambda=
\begin{pmatrix}
-1\\
-\lambda
\end{pmatrix}.
\]
В работе используются начальные значения
\[
\lambda_0\in\{-2,-1.5,-1\}.
\]

\paragraph{case5: устойчивость при аффинном терминальном ограничении}
Используется та же двумерная система, что и в case4,
\[
A=
\begin{pmatrix}
0&1\\
-1&0
\end{pmatrix},
\qquad
B=
\begin{pmatrix}
0\\
1
\end{pmatrix},
\qquad |u|\leq 1.
\]
но с другими $x_0$ и $d$:
\[
x_0=
\begin{pmatrix}
1\\
0
\end{pmatrix},
\quad
T=2,
\quad
a=
\begin{pmatrix}
1\\
0
\end{pmatrix}.
\]
Терминальное множество задаётся как
\[
M=\{x\in\mathbb{R}^2:Cx=d\},
\quad
C=
\begin{pmatrix}
0&1
\end{pmatrix},
\quad
d=0.
\]
Здесь также
\[
q(\lambda)=-a-C^\top\lambda=
\begin{pmatrix}
-1\\
-\lambda
\end{pmatrix},
\]
а терминальная невязка ограничения определяется формулой
\[
G(\lambda)=Cx_M(\lambda)-d.
\]
Начальные приближения выбираются как
\[
\lambda_0\in\{-3,-2.5,-2\}.
\]

\subsection{5.5 Сводка результатов и обсуждение}

В этом разделе кратко излагаются ключевые результаты case1--case5 и их доказательная роль.

\subsubsection*{case1: метод, точность и сетка в двумерной задаче терминального отклонения}
\begin{figure}[htbp]
\centering
\includegraphics[width=0.94\textwidth]{code_v1.4/results/figures/case1_panel.png}
\caption{Обзор результатов case1: геометрия терминального отклонения, базовый уровень точности и влияние сетки}
\label{fig:case1}
\end{figure}

На рис.~\ref{fig:case1} показано сравнение различных реализаций первой задачи. Метод Гаусса--Ньютона достигает критерия остановки за 5 итераций, при этом невязка составляет около $2.05\times 10^{-8}$, а терминальная ошибка --- около $1.49\times 10^{-3}$. Градиентный метод не достигает сходимости за 120 итераций, сохраняя невязку порядка $3.15\times 10^{-1}$ и терминальную ошибку порядка $2.90\times 10^{-1}$. При той же сетке $M=200$ прямой дискретный метод даёт терминальную ошибку около $3.68\times 10^{-9}$.

Этот пример поддерживает три осторожных вывода: в текущей реализации направление Гаусса--Ньютона даёт существенно более быстрое убывание невязки, чем градиентное направление; прямой метод на текущей сетке и при текущем критерии остановки обеспечивает меньшую терминальную ошибку; изменение числа шагов от $50$ до $800$ меняет терминальную ошибку метода Гаусса--Ньютона лишь в узком диапазоне от примерно $1.46\times 10^{-3}$ до $1.50\times 10^{-3}$. Поэтому case1 следует использовать как реализационный ориентир и точностной базис, а не как доказательство сильного закона сеточной сходимости.

\subsubsection*{case2: как параметр сглаживания одновременно влияет на сходимость и форму управления}
Для case2 при $\mu=10^{-1},10^{-2},10^{-3}$ метод Гаусса--Ньютона сходится за 4--5 итераций, а терминальная ошибка составляет соответственно примерно $1.09\times 10^{-1}$, $1.43\times 10^{-2}$ и $1.49\times 10^{-3}$. При $\mu=0$ алгоритм не сходится за 120 итераций, а и невязка, и терминальная ошибка остаются порядка $2.95\times 10^{-1}$. Этот пример подтверждает, что сглаживание влияет одновременно на устойчивость сходимости и на близость к bang-bang структуре.

\subsubsection*{case3: меньшая размерность переменной и граница точности в четырёхмерной системе}
\begin{figure}[htbp]
\centering
\includegraphics[width=0.94\textwidth]{code_v1.4/results/figures/case3_panel.png}
\caption{Обзор результатов case3: уменьшение размерности переменной и сравнение терминальной точности}
\label{fig:case3}
\end{figure}

На рис.~\ref{fig:case3} показано, что низкая размерность переменной и высокая точность не обязательно достигаются одновременно. В текущей реализации и при выбранных критериях остановки сопряжённый метод Гаусса--Ньютона достигает остаточного критерия за 4 итерации, тогда как прямой дискретный метод использует 8 итераций. При этом косвенный метод оптимизирует только 4 параметра, а прямой --- 400 дискретных переменных управления. Однако терминальная ошибка косвенного метода составляет около $1.08\times 10^{-3}$, тогда как прямой метод даёт уровень примерно $1.17\times 10^{-9}$.

Следовательно, основная ценность параметризации по сопряжённой переменной заключается в уменьшении размерности и в сохранении структурной интерпретации, а не в гарантии лучшей терминальной точности во всех задачах средней размерности. В этом примере прямой метод остаётся важным точностным ориентиром.

\subsubsection*{case4: геометрия терминального множества и линейная терминальная цель}
В case4 для начальных значений $\lambda_0=-2,-1.5,-1$ число итераций составляет 8, 7 и 10, невязка терминального ограничения имеет порядок $10^{-8}$--$10^{-7}$, а итоговое значение терминальной цели стабилизируется около $-1.53289\times 10^{-1}$. Этот пример показывает, что параметризация по множителю приводит к близким допустимым терминальным состояниям на целевом множестве и тем самым поддерживает геометрическую интерпретацию второй задачи.

\subsubsection*{case5: выполнение терминального ограничения и устойчивость по начальному множителю}
\begin{figure}[htbp]
\centering
\includegraphics[width=0.94\textwidth]{code_v1.4/results/figures/case5_panel.png}
\caption{Обзор результатов case5: устойчивость при разных начальных значениях множителя}
\label{fig:case5}
\end{figure}

На рис.~\ref{fig:case5} показана устойчивость второй задачи по начальному приближению множителя. Для начальных значений $\lambda_0=-3,-2.5,-2$ число итераций равно 7, 9 и 7, невязка терминального ограничения имеет порядок $1.73\times 10^{-7}$, $3.93\times 10^{-9}$ и $7.97\times 10^{-7}$, а итоговое значение цели составляет примерно $-0.6602688$, $-0.6602692$ и $-0.6602674$.

Из этого следует, что для второй задачи более информативно отслеживать устойчивость невязки терминального ограничения и итогового значения линейной цели, чем говорить о ``терминальной ошибке'' по аналогии с первой задачей. В текущем примере изменение начального множителя влияет на промежуточную итерационную траекторию, но не меняет конечный уровень допустимости и итоговую цель.

\paragraph{Итоговое замечание}
Пять примеров совместно показывают, что низкоразмерная параметризация вычислима и интерпретируема для обоих классов задач. Одновременно они подтверждают, что при негладкости, вырождении и повышенных требованиях к точности прямой метод остаётся важным ориентиром.

\section{Выводы}

Основные выводы работы таковы:
\begin{enumerate}
\item Для линейных систем с постоянными коэффициентами сопряжённая траектория допускает низкоразмерную параметризацию по терминальному вектору или множителю.
\item Построенные уравнения невязки представляют собой дискретную реализацию необходимых условий принципа максимума и служат средством построения кандидатных решений.
\item Для первой задачи прямой метод остаётся важной базой точности, а для второй ключевыми интерпретационными величинами являются невязка ограничения и устойчивость итогового значения функционала.
\item При вырождении, негладкости и более жёстких требованиях к точности предложенный подход естественно дополняется прямым методом, регуляризацией или дополнительными правилами выбора.
\end{enumerate}

\section{Глава 6. Заключение}

Позицию предложенного подхода можно кратко сформулировать в трёх пунктах. Во-первых, он сжимает две линейные терминальные задачи до более низкоразмерных задач поиска невязки, сохраняя структурную связь с принципом максимума и терминальной параметризацией. Во-вторых, он не даёт оснований напрямую утверждать, что вычисляется строгий оптимум непрерывной исходной задачи или явной Euler-дискретизации; корректнее рассматривать результат как дискретное кандидатное решение, согласованное с непрерывными необходимыми условиями. В-третьих, прямой дискретный метод сохраняет незаменимую роль в качестве точностного ориентира, инструмента проверки достижимости и средства анализа вырожденных режимов.

Перспективы дальнейшей работы естественно связаны с более строгим анализом дискретной сходимости, с расширением класса допустимых терминальных ограничений и с построением более устойчивых численных процедур для негладких и вырожденных режимов.

\section{Список литературы}

{\sloppy\raggedright
\begin{thebibliography}{99}
\bibitem{pontryagin1969}
Л. С. Понтрягин, В. Г. Болтянский, Р. В. Гамкрелидзе, Е. Ф. Мищенко.
\newblock \emph{Математическая теория оптимальных процессов}.
\newblock Москва: Наука, 1969.

\bibitem{kiselev2007}
Ю. Н. Киселёв, С. Н. Аввакумов, М. В. Орлов.
\newblock \emph{Оптимальное управление. Линейная теория и приложения}.
\newblock Москва: Издательский отдел факультета ВМиК МГУ им. М. В. Ломоносова, 2007.

\bibitem{lee1972}
Э. Б. Ли, Л. Маркус.
\newblock \emph{Основы теории оптимального управления}.
\newblock Москва: Наука, 1972.

\bibitem{blagodatskikh2001}
В. И. Благодатских.
\newblock \emph{Введение в оптимальное управление (линейная теория)}.
\newblock Москва: Высшая школа, 2001.

\bibitem{pontryagin1985}
Л. С. Понтрягин.
\newblock Математическая теория оптимальных процессов и дифференциальные игры.
\newblock \emph{Труды Математического института АН СССР}, 1985, т. 169, с. 119--158.

\bibitem{avvakumov1995}
С. Н. Аввакумов, Ю. Н. Киселёв, М. В. Орлов.
\newblock Методы решения задач оптимального управления на основе принципа максимума Понтрягина.
\newblock \emph{Труды Математического института им. В. А. Стеклова}, 1995, т. 211, с. 3--31.

\bibitem{bryson1975}
A. E. Bryson, Y.-C. Ho.
\newblock \emph{Applied Optimal Control: Optimization, Estimation, and Control}.
\newblock Washington: Hemisphere Publishing, 1975.

\bibitem{kirk1970}
D. E. Kirk.
\newblock \emph{Optimal Control Theory: An Introduction}.
\newblock Englewood Cliffs: Prentice-Hall, 1970.

\bibitem{vinter2000}
R. B. Vinter.
\newblock \emph{Optimal Control}.
\newblock Boston: Birkhäuser, 2000.

\bibitem{liberzon2012}
D. Liberzon.
\newblock \emph{Calculus of Variations and Optimal Control Theory: A Concise Introduction}.
\newblock Princeton: Princeton University Press, 2012.

\bibitem{betts2010}
J. T. Betts.
\newblock \emph{Practical Methods for Optimal Control and Estimation Using Nonlinear Programming}.
\newblock Philadelphia: SIAM, 2010.

\bibitem{rao2009}
A. V. Rao.
\newblock A survey of numerical methods for optimal control.
\newblock \emph{Advances in the Astronautical Sciences}, 2009, vol. 135, pp. 497--528.

\bibitem{nocedal2006}
J. Nocedal, S. J. Wright.
\newblock \emph{Numerical Optimization}. 2nd ed.
\newblock New York: Springer, 2006.

\bibitem{kelley1999}
C. T. Kelley.
\newblock \emph{Iterative Methods for Optimization}.
\newblock Philadelphia: SIAM, 1999.

\bibitem{bock1984}
H. G. Bock, K.-J. Plitt.
\newblock A multiple shooting algorithm for direct solution of optimal control problems.
\newblock In: \emph{Proceedings of the 9th IFAC World Congress}, Budapest, 1984, pp. 242--247.

\bibitem{hager2000}
W. W. Hager.
\newblock Runge--Kutta methods in optimal control and the transformed adjoint system.
\newblock \emph{Numerische Mathematik}, 2000, vol. 87, no. 2, pp. 247--282.

\bibitem{dontchev2000}
A. L. Dontchev, W. W. Hager, V. M. Veliov.
\newblock Second-order Runge--Kutta approximations in control constrained optimal control.
\newblock \emph{SIAM Journal on Numerical Analysis}, 2000, vol. 38, no. 1, pp. 202--226.

\bibitem{maurer1995}
H. Maurer, H. J. Pesch.
\newblock Solution differentiability for nonlinear parametric control problems.
\newblock \emph{SIAM Journal on Control and Optimization}, 1995, vol. 33, no. 5, pp. 1456--1483.

\bibitem{bonnans2009}
J. F. Bonnans, A. Hermant.
\newblock Second-order analysis for optimal control problems with pure state constraints and mixed control-state constraints.
\newblock \emph{Annales de l'Institut Henri Poincare C}, 2009, vol. 26, no. 2, pp. 561--598.

\bibitem{github-code}
Сяо Ижу.
\newblock \emph{Numerical solution of the terminal optimal control problem with linear dynamics}.
\newblock Репозиторий GitHub.
\newblock \href{https://github.com/yxforever666gh/Numerical-solution-of-the-terminal-optimal-control-problem-with-linear-dynamics.git}{Repository link}.
\newblock Дата обращения: 2026-05-04.
\end{thebibliography}
\fussy
}

\end{document}
